\documentclass[11pt]{article}
\usepackage[margin=1in]{geometry}
\usepackage{amsmath,amssymb,amsthm}
\usepackage{booktabs}
\usepackage{enumitem}
\usepackage{graphicx}
\usepackage{hyperref}

\newcommand{\dd}{\mathrm{d}}
\newcommand{\norm}[1]{\lVert #1 \rVert}
\newcommand{\inner}[2]{\left( #1 , #2 \right)}
\newcommand{\code}[1]{\texttt{#1}}
\newcommand{\bs}[1]{\boldsymbol{#1}}
\newtheorem{proposition}{Proposition}
\newtheorem{remark}{Remark}
\emergencystretch=3em

\title{Variational Recovery of Internal Variables\\
in Admissible Spaces:\\
Constrained $L^2$ Projection for Norm-Like Variables}
\author{Alejandro Mota}
\date{August 2026}

\begin{document}
\maketitle

\begin{abstract}
The variational recovery of internal variables of
Mota~et~al.~\cite{mota2013} extends integration-point fields to the entire
domain by an $L^2$ projection derived from a three-field variational
principle, and preserves the admissible spaces of group-valued variables by
interpolating in the corresponding Lie algebras through the logarithmic and
exponential maps. This treatment leaves open the case of norm-like
variables such as the equivalent plastic strain, whose admissible set is
the closed half-line $[0,\infty)$: the field attains the boundary value
zero over finite regions, the half-line is not a group, and in practice the
consistent projection produces negative values near steep gradients. This
note argues that the correct extension is not a repair of the logarithmic
transformation but a change of the admissible set in the same variational
principle: minimize the same $L^2$ distance over the intersection of the
finite element space with the convex cone of admissible fields. The
resulting constrained projection is the unique solution of a variational
inequality; it is idempotent and non-expansive by construction, properties
that the unconstrained projection followed by clamping does not possess.
Discretely it is a bound-constrained quadratic program on the
already-assembled Gram matrix, solvable by a primal--dual active-set method
in a small, finite number of iterations. The active set of the constraint
is a localized indicator of an under-resolved front and therefore serves
directly as an input to error estimation and mesh size-field
construction. The note
closes with a taxonomy of internal variables by admissible set and an
outline of an implementation in Norma.
\end{abstract}

\tableofcontents

\section{Background and problem statement}
\label{sec:background}

The recovery procedure of \cite{mota2013} arises from a three-field
functional in the deformation mapping $\bs{\varphi}$, a
\emph{target} internal variable field $\bar{\bs{z}}$ defined over the whole
domain, and a Lagrange multiplier $\bar{\bs{y}}$ that enforces equality
between the target field and the \emph{source} field $\bs{z}$ available at
integration points. Discretization with interpolation functions
$\lambda_\alpha$ for the target field yields the consistent $L^2$
projection
\begin{equation}
\bar{\bs{z}}_h(\bs{X}) = \lambda_\alpha(\bs{X})\,\bar{\bs{z}}_\alpha,
\qquad
\bar{\bs{z}}_\alpha = \bs{H}_{\alpha\beta}^{-1}\,\bs{\Sigma}_\beta,
\qquad
\bs{H}_{\alpha\beta} := \int_B \lambda_\alpha \lambda_\beta\, \bs{I} \, \dd V,
\qquad
\bs{\Sigma}_\beta := \int_B \lambda_\beta\, \bs{z} \, \dd V,
\label{eq:projection}
\end{equation}
with $\bs{I}$ the $q \times q$ identity and $q$ the number of scalar
parameters that represent the internal variable. The projection
\eqref{eq:projection} is the minimizer of
$\tfrac12 \norm{\bar{\bs{z}}_h - \bs{z}}_{L^2(B)}^2$ over the finite
element space $V_h$. The Gram matrix $\bs{H}$ is symmetric positive
definite, sparse, and of narrow band. Because the projection is a
minimizer, it is optimal; because it is a projection onto a subspace, it is
idempotent, which is the property that makes it suitable for repeated
remap: lumped (diagonalized) variants are not idempotent and diffuse the
field under repeated application.

For internal variables that do not live in a linear space, \cite{mota2013}
interpolates in the Lie algebra: the field $\bs{Z}$ with values in
a Lie group $G$ is transformed by the logarithmic map to $\bs{z} = \log
\bs{Z}$ in the Lie algebra $\mathfrak{g}$, which is a vector space;
the projection \eqref{eq:projection} is applied there; and the result is
returned to the group by the exponential map. This guarantees, for
example, that projected rotations remain rotations, that isochoric
deformations remain isochoric, and that a volumetric Jacobian remains
positive.

The open question addressed here is the treatment of \emph{norm-like}
variables, of which the equivalent plastic strain
$\varepsilon^p = \int \dot{\varepsilon}^p \,\dd t \geq 0$
is the canonical example. Its admissible set is the closed half-line
$[0,\infty)$. In practice, the consistent projection of such a variable
produces negative values near steep gradients --- the projection is a
(near-)best $L^2$ approximation, and best approximations of fields with
sharp fronts oscillate. Negative equivalent plastic strain subsequently
enters constitutive updates that assume non-negativity, with consequences
ranging from silent inaccuracy to failure of the local return mapping.

\section{Why the logarithmic transformation does not extend to this case}
\label{sec:limits}

It is tempting to treat $[0,\infty)$ by analogy with the positive reals
$\mathbb{R}^+$, which form a Lie group under multiplication with Lie
algebra $\mathbb{R}$. The analogy fails for a structural reason, and
identifying that reason yields a criterion that predicts which variables
need which treatment.

The half-line $[0,\infty)$ is not a group: zero has no multiplicative
inverse, and $\log 0 = -\infty$. More importantly, zero is not an
exceptional value of the equivalent plastic strain but its \emph{generic}
value: every elastic region carries $\varepsilon^p = 0$ exactly, over
finite subdomains. The field therefore sits \emph{on} the boundary of its
admissible set, not merely near it.

Contrast this with the variables the logarithmic transformation serves
well. The volumetric Jacobian $J = \det \bs{F}$ is positive, and
$J \to 0$ is approached only asymptotically under extreme compression;
the generic state $J = 1$ maps to $\log J = 0$, the identity of the
algebra. Scalar damage $D \in [0,1)$, transformed as
$1 - D \in \mathbb{R}^+$ with $h = \log(1 - D)$, maps the generic
virgin state $D = 0$ to $h = 0$; the boundary $D \to 1$ is asymptotic
and is pushed to $h \to -\infty$, where linear interpolation cannot
overshoot it.

The criterion is then:
\begin{quote}
\emph{Asymptotic boundaries --- approached but never attained --- are
handled by a logarithmic transformation, which maps them to infinity.
Attained boundaries --- on which the field sits over finite regions ---
cannot be transformed away and must be enforced as inequality
constraints.}
\end{quote}
The two mechanisms compose. Scalar damage possesses both kinds of
boundary: the asymptotic boundary $D \to 1$ is removed by the
transformation $h = \log(1-D) \in (-\infty, 0]$, while the attained
boundary $D = 0$, that is $h = 0$, persists in the transformed variable
and requires the one-sided constraint $h \leq 0$ --- which is again a
convex cone constraint, now linear in the transformed space. The
equivalent plastic strain has only the attained boundary and is treated
by a cone constraint in physical space directly.

A regularized transformation $\log(\epsilon + \varepsilon^p)$ is not a
remedy. First, the reconstructed location of an elastic--plastic front
becomes a function of the regularization parameter $\epsilon$ rather than
of the data. Second, the $L^2$ norm in the transformed space weights
\emph{relative} error: fluctuations of order $\log \epsilon$ in the
elastic region contribute to the minimized distance as strongly as the
physically meaningful variation in the plastic zone, and the minimizer
approximates the former in preference to the latter.

\section{Constrained projection as a variational inequality}
\label{sec:vi}

The proposed extension retains the variational principle of
\cite{mota2013} verbatim and changes only the admissible set. Let
$K \subseteq \mathbb{R}^q$ be the closed convex set of admissible values
of the internal variable (a cone, a box, or a simplex; see
Section~\ref{sec:taxonomy}), and define the admissible subset of the
finite element space
\begin{equation}
\mathcal{K}_h := \left\{ \bs{\eta}_h \in V_h \;:\;
\bs{\eta}_h(\bs{X}) \in K \ \text{for all } \bs{X} \in B
\right\},
\end{equation}
which is closed and convex. The constrained recovery is
\begin{equation}
\bar{\bs{z}}_h = \operatorname*{arg\,min}_{\bs{\eta}_h \in \mathcal{K}_h}
\tfrac12 \norm{\bs{\eta}_h - \bs{z}}_{L^2(B)}^2 ,
\label{eq:constrained}
\end{equation}
equivalently the solution of the variational inequality
\begin{equation}
\bar{\bs{z}}_h \in \mathcal{K}_h :
\qquad
\inner{\bar{\bs{z}}_h - \bs{z}}{\bs{\eta}_h - \bar{\bs{z}}_h}_{L^2(B)} \geq 0
\qquad \forall\, \bs{\eta}_h \in \mathcal{K}_h ,
\label{eq:vi}
\end{equation}
the discrete analogue of the obstacle problem. When no constraint is
active, \eqref{eq:constrained} reduces identically to
\eqref{eq:projection}; on a mesh that resolves the field, the constrained
and unconstrained recoveries coincide.

The essential properties follow from the Hilbert space projection theorem
onto closed convex sets.

\begin{proposition}[Well-posedness, idempotency, stability]
\label{prop:properties}
Problem \eqref{eq:constrained} admits a unique solution
$\bar{\bs{z}}_h = P_{\mathcal{K}_h}\, \bs{z}$. The map $P_{\mathcal{K}_h}$
is idempotent, $P_{\mathcal{K}_h} \circ P_{\mathcal{K}_h} =
P_{\mathcal{K}_h}$, and non-expansive,
$\norm{P_{\mathcal{K}_h} \bs{u} - P_{\mathcal{K}_h} \bs{v}} \leq
\norm{\bs{u} - \bs{v}}$.
\end{proposition}

Idempotency deserves emphasis because it has been raised as a concern:
one might fear that appending constraints to the projection destroys the
idempotency that motivates the consistent projection in the first place.
The opposite holds. The metric projection onto a closed convex set is
idempotent \emph{by construction} --- an admissible field is its own
closest admissible point --- whereas the apparently simpler alternative,
the unconstrained projection followed by a pointwise clamp, is not: the
clamp is the projection in the wrong (diagonal) metric, the clamped field
is generally not the $L^2$-closest admissible field, and re-projecting it
re-introduces oscillation. Under the repeated project--interpolate cycles
of successive remaps, the composition ``project, then clamp'' drifts; the
constrained projection does not. Non-expansiveness supplies the
complementary stability statement: a sequence of constrained projections
cannot amplify perturbations.

\begin{remark}[One-sided conservation of the accumulated quantity]
\label{rem:mean}
Let $K = [0,\infty)^q$, componentwise non-negativity of the $q$ scalar
parameters that represent the internal variable in
\eqref{eq:projection}, with the equivalent plastic strain as the scalar
case $q = 1$. The \emph{unconstrained} projection conserves the integral
of the field: the Lagrange interpolation functions form a partition of
unity, $\sum_\alpha \lambda_\alpha = 1$, so constant fields belong to
$V_h$, and testing the unconstrained stationarity condition
$\inner{\bar{\bs{z}}_h - \bs{z}}{\bs{\eta}_h}_{L^2(B)} = 0$ with
$\bs{\eta}_h = \bs{1}$, the constant field of ones in $\mathbb{R}^q$,
yields
$\int_B \bs{1} \cdot \bar{\bs{z}}_h \,\dd V =
\int_B \bs{1} \cdot \bs{z} \,\dd V$.

The constrained projection retains exactly one half of this statement,
and the geometry of the cone selects which half. In the variational
inequality \eqref{eq:vi} the perturbation
$\bs{\eta}_h = \bar{\bs{z}}_h + c\,\bs{1}$, with $c > 0$ a scalar
constant, is always admissible --- raising an admissible field keeps it
admissible --- and substitution gives
$c \int_B (\bar{\bs{z}}_h - \bs{z}) \cdot \bs{1} \,\dd V \geq 0$,
that is,
\begin{equation}
\int_B \bs{1} \cdot \bar{\bs{z}}_h \,\dd V \;\geq\;
\int_B \bs{1} \cdot \bs{z} \,\dd V .
\label{eq:mean}
\end{equation}
The opposite perturbation $\bs{\eta}_h = \bar{\bs{z}}_h - c\,\bs{1}$
is inadmissible wherever any component of the solution touches the
bound, so no matching upper inequality survives: the projection may
raise the total of the field, never lower it.

Discretely the statement sharpens to an identity. In the
scalar-variable setting of Section~\ref{sec:discrete}, multiplying the
complementarity system \eqref{eq:kkt} by $\bs{1}^{\mathsf{T}}$, with
$\bs{1}$ the vector of ones, and using the partition of unity twice,
namely
$\bs{1}^{\mathsf{T}} \bs{H} \bar{\bs{z}} = \int_B \bar{\bs{z}}_h \,\dd V$
and $\bs{1}^{\mathsf{T}} \bs{\Sigma} = \int_B \bs{z} \,\dd V$, gives
\begin{equation}
\int_B \bar{\bs{z}}_h \,\dd V - \int_B \bs{z} \,\dd V
= \bs{1}^{\mathsf{T}} \bs{\mu} \;\geq\; 0 ,
\label{eq:meanid}
\end{equation}
so the excess of the accumulated variable equals the resultant of the
admissibility reactions $\bs{\mu}$ of
Remark~\ref{rem:multiplier}. It is therefore supported on the active
set, vanishes identically when no constraint activates, and decays with
the local projection error as the mesh resolves the front.

The direction of inequality \eqref{eq:mean} is the physically correct
one. Repairing a negative undershoot of the equivalent plastic strain
must add plastic history in a neighborhood of the front; a repair able
to remove it would fabricate less-yielded material out of a numerical
oscillation. Among all admissible fields the constrained projection adds
the least, since $\bar{\bs{z}}_h$ is by definition the closest
admissible field to the source.
\end{remark}

\begin{remark}[Interpretation of the multiplier]
\label{rem:multiplier}
In \cite{mota2013} the Lagrange multiplier of the three-field functional
is identified with the thermodynamic force conjugate to the internal
variable. Appending the inequality $\bar{\bs{z}}_h \in \mathcal{K}_h$
introduces a further multiplier $\bs{\mu} \geq \bs{0}$, complementary to
the constraint, which plays the role of the reaction that maintains
admissibility of the state --- the same role the contact pressure plays
for the impenetrability constraint. Its support marks precisely the
regions where the unconstrained projection would exit the admissible set;
Section~\ref{sec:adaptivity} exploits this.
\end{remark}

\section{Discrete formulation and solution algorithms}
\label{sec:discrete}

\subsection{Bound-constrained quadratic program}

For a scalar variable with $K = [0,\infty)$ and constraints imposed on
the vector $\bar{\bs{z}}$ of nodal values (the adequacy of nodal
constraints is discussed in Section~\ref{sec:where}), problem
\eqref{eq:constrained} is the bound-constrained quadratic program
\begin{equation}
\min_{\bar{\bs{z}} \geq \bs{0}} \
\tfrac12\, \bar{\bs{z}}^{\mathsf{T}} \bs{H} \bar{\bs{z}}
- \bs{\Sigma}^{\mathsf{T}} \bar{\bs{z}},
\label{eq:qp}
\end{equation}
with $\bs{H}$ and $\bs{\Sigma}$ exactly as assembled for
\eqref{eq:projection}. The first-order conditions are the linear
complementarity system
\begin{equation}
\bs{H} \bar{\bs{z}} - \bs{\Sigma} = \bs{\mu}, \qquad
\bar{\bs{z}} \geq \bs{0}, \qquad \bs{\mu} \geq \bs{0}, \qquad
\bs{\mu}^{\mathsf{T}} \bar{\bs{z}} = 0 .
\label{eq:kkt}
\end{equation}
Because $\bs{H}$ is symmetric positive definite, \eqref{eq:qp} is
strictly convex and possesses exactly one solution: the closest
admissible field of \eqref{eq:constrained}. The two algorithms below
compute this same point and differ only in how they reach it --- the
first exactly, in finitely many steps; the second in the limit of an
infinite iteration.

\subsection{Primal--dual active set (preferred for a new implementation)}
\label{sec:pdas}

The primal--dual active-set method --- equivalent to a semismooth Newton
method on the complementarity system \cite{hintermuller2003} ---
terminates at the exact solution of \eqref{eq:qp} in a finite number of
iterations, typically fewer than five in practice. Each iteration fixes the currently
active nodes at the bound and solves the reduced symmetric positive
definite system on the inactive set. A factorization of $\bs{H}$,
computed once per remap, is reused across iterations, so the additional
computation relative to the unconstrained projection
\eqref{eq:projection} amounts to a small number of back-substitutions
on subsets of the unknowns. When
no constraints activate, the method terminates after one iteration with
the unconstrained solution.

Two further properties recommend it in the remap setting. First, it
warm-starts: the front moves little between successive remaps, so the
active set of the previous remap is an excellent initial guess, and the
method typically confirms or barely amends it in one or two iterations.
Second, it yields the multiplier $\bs{\mu}$ exactly at termination,
and the multiplier is precisely the quantity that
Section~\ref{sec:adaptivity} supplies to the size field: a truncated
iterative method produces $\bs{\mu}$ only to the accuracy of its
stopping tolerance, contaminating the small, localized indicator one
wishes to read.

\subsection{Dykstra's alternating projections (preferred when
retrofitting a legacy code)}
\label{sec:dykstra}

An implementation that prefers to compose existing operators rather than
introduce a quadratic programming kernel may use Dykstra's algorithm
\cite{boyle1986}: alternate the consistent projection onto $V_h$ (a Gram
matrix solve) and the pointwise clamp onto the cone of non-negative
fields, with Dykstra's correction increments carried between passes.
The correction increments are essential, and their role is the point
most easily misread. \emph{Plain} alternating projections --- project,
clamp, repeat, with no corrections --- converge to \emph{some} point of
the intersection of $V_h$ with the cone: an admissible field, but in
general not the minimizer of \eqref{eq:qp}, hence a suboptimal recovery.
Dykstra's increments repair exactly this defect: the corrected iteration
converges to the solution of \eqref{eq:qp}, the same point the
active-set method computes. The difference between the two subsections
is therefore not the answer but the route: the active-set method reaches
it exactly in finitely many steps, whereas Dykstra's iteration reaches
it only in the limit, at a linear rate that degrades precisely where the
constraint matters --- near a sharp front, where the subspace and the
cone meet at a shallow angle.

The advantage of Dykstra's algorithm is therefore architectural rather
than mathematical: it introduces no new solver kernel, only the
composition of two operators a legacy code already possesses --- the
consistent Gram matrix solve and the pointwise clamp. In an existing
code where a bound-constrained solver would be an intrusive addition,
it is the appropriate first implementation; in a new implementation,
where the active-set solver of Section~\ref{sec:pdas} is
straightforward to write, nothing else favors it.

\subsection{Quadrature consistency}

The unconstrained projection is idempotent under inexact quadrature
provided $\bs{H}$ and $\bs{\Sigma}$ are integrated with the same rule:
the projection of a member of $V_h$ then reproduces it identically. The
constrained projection inherits this property, since an admissible member
of $V_h$ satisfies \eqref{eq:vi} trivially. The discipline of
\cite{mota2013} --- the same integration points for the recovery as for
the equilibrium statement --- carries over unchanged.

\section{Where to impose the bounds}
\label{sec:where}

Nodal constraints $\bar{\bs{z}}_\alpha \geq \bs{0}$ guarantee a
non-negative \emph{field} only when the interpolation functions are
non-negative, as for linear simplices and multilinear hexahedra. For
quadratic and serendipity elements the basis functions change sign and
nodal bounds do not imply field bounds. Two options exist:
\begin{enumerate}[nosep]
\item constrain the field at the points where it will be evaluated ---
the integration points of the target mesh --- which is natural for
in-core remap where the target is known at projection time; or
\item express the element polynomial in the Bernstein basis and constrain
the Bernstein coefficients, whose non-negativity is a sufficient (mildly
conservative) condition for non-negativity of the polynomial. This is the
standard device of bound-preserving high-order transport and requires no
knowledge of the target mesh, which suits file-based remap.
\end{enumerate}
For the linear elements that dominate current practice, nodal constraints
suffice.

\section{A taxonomy of internal variables}
\label{sec:taxonomy}

Table~\ref{tab:taxonomy} classifies common internal variables by
admissible set, boundary character, and the resulting treatment. The
classification is the operational content of the criterion of
Section~\ref{sec:limits}.

\begin{table}[htbp]
\centering
\small
\begin{tabular}{llll}
\toprule
Variable & Admissible set & Boundary character & Treatment \\
\midrule
Back stress, strains & vector space & none & projection \eqref{eq:projection} \\
Rotations $\bs{R}$ & $SO(3)$ & none (group) & $\log$/$\exp$ + \eqref{eq:projection} \\
Isochoric $\bs{F}^p$ & $SL(3)$ & none (group) & $\log$/$\exp$ + \eqref{eq:projection} \\
Jacobian $J$ & $(0,\infty)$ & $0$ asymptotic & $\log$ + \eqref{eq:projection} \\
Equiv.\ plastic strain $\varepsilon^p$ & $[0,\infty)$ & $0$ attained & cone constraint \eqref{eq:constrained} \\
Scalar damage $D$, porosity $f$ & $[0,1)$ & $0$ attained, $1$ asymptotic &
$h = \log(1-D)$, constraint $h \leq 0$ \\
Volume fractions $\phi_i$ & simplex & faces attained &
\eqref{eq:constrained} with $\sum_i \phi_i = 1$ \\
\bottomrule
\end{tabular}
\caption{Internal variables by admissible set. ``Attained'' boundaries
are those on which the field sits over finite regions; they require
inequality constraints. ``Asymptotic'' boundaries are removed by
logarithmic transformation. The two mechanisms compose, as in the damage
row.}
\label{tab:taxonomy}
\end{table}

The architectural consequence is that the admissible set is
\emph{metadata of the state variable}, declared by the constitutive model
and read by the transfer machinery, which dispatches on it: plain
projection for vector-space variables, logarithmic transformation for
group-valued variables, constrained projection for cone- and box-valued
variables, and their composition where the taxonomy requires it. Codes
that already tag variables for transfer (for instance, as eligible for
``$L^2$ projection'') possess the natural place for this declaration; the
tag needs only to be enriched from a method label to an admissible-set
descriptor.

A complementary safeguard belongs in the constitutive model itself:
state entering a material update --- from initial conditions, restart, or
transfer --- should be checked against the declared admissible set, with
tolerance semantics: violations below a tolerance are repaired silently,
violations above it are reported as errors. This is a verification
layer, not the primary mechanism. A silent clamp buried in the material
model conceals transfer
defects and re-introduces exactly the non-idempotent repair that
Section~\ref{sec:vi} argues against; if the transfer enforces
admissibility variationally, the material-level check should never
activate.

\section{Interaction with error estimation and mesh adaptation}
\label{sec:adaptivity}

The negativity that motivates this note is a symptom; the underlying
cause is a front that the mesh does not resolve. The constrained projection removes the
inadmissibility but not the approximation error, and it is therefore a
complement to, not a substitute for, adaptive refinement driven by the
recovery error. The two interact favorably in both directions.

First, the multiplier $\bs{\mu}$ of \eqref{eq:kkt}, obtained as a
byproduct of the solution, is a localized indicator of exactly the
regions where the finite element space cannot represent the front: its support is the active set, which is nonempty
precisely where the unconstrained projection would exit the admissible
set. In a remap-driven adaptive cycle --- project the field, interpolate
back to the integration points, measure the $L^2$ discrepancy elementwise,
construct a size field, remesh --- the multiplier supplies an additional
element-level input to the size field without further computation,
identifying fronts that the plain $L^2$ discrepancy may under-weight
because the offending oscillations are small in norm while large in
consequence.

Second, as refinement resolves the
front, the constraints deactivate and the constrained projection reduces
continuously to the unconstrained one, whose optimality and idempotency
are already established. The constraint governs the pre-asymptotic
regime; refinement removes the need for it. Near a genuine free boundary
the recovered field remains only $C^0$ at any resolution, so mild
residual constraint activity at the front is expected and benign.

\section{Approaches considered and rejected}
\label{sec:rejected}

\begin{description}[nosep]
\item[Clamping after projection.] Not idempotent; the projection in the
wrong metric; drifts under repeated remap
(Section~\ref{sec:vi}).
\item[Regularized logarithm $\log(\epsilon + \varepsilon^p)$.] Front
location depends on $\epsilon$; relative-error weighting resolves the
elastic region in preference to the plastic zone
(Section~\ref{sec:limits}).
\item[Lumped projection.] Idempotency fails; repeated remap diffuses the
field. This is the original motivation for the consistent projection of
\cite{mota2013} and is not revisited here.
\item[Conservation constraint on the mean.] Appending
$\int_B \bar{\bs{z}}_h \,\dd V = \int_B \bs{z} \,\dd V$ to
$\mathcal{K}_h$ is convex and implementable, but the equivalent plastic
strain is not a conserved quantity; the one-sided inequality
\eqref{eq:mean} is the appropriate statement, and forcing equality
would spread the front's admissibility correction over the domain.
\end{description}

\section{Outline of an implementation in Norma}
\label{sec:norma}

Norma is well suited for a prototype: it already assembles consistent
mass (Gram) matrices, carries constitutive state at integration points,
and possesses the quadrature and interpolation infrastructure of
\cite{mota2013}. A minimal introduction would comprise:
\begin{enumerate}[nosep]
\item a recovery module implementing \eqref{eq:projection} with the Gram
matrix assembled by the existing interpolation and quadrature routines,
verified for idempotency under project--interpolate cycles on a fixed
mesh;
\item the admissible-set tag on constitutive state
(vector space, Lie group with its $\log$/$\exp$ pair, cone, box,
simplex), with the dispatch of Section~\ref{sec:taxonomy};
\item the primal--dual active-set solver for \eqref{eq:qp}, exercised
first on a manufactured front (a piecewise field with an attained-boundary
region) and then on a Taylor-bar-type problem under repeated remap, with
the clamped unconstrained projection as the control demonstrating drift;
\item the multiplier-to-size-field hook of
Section~\ref{sec:adaptivity}, deferred until a remeshing loop exists.
\end{enumerate}
Items 1--3 are self-contained and testable without any remeshing
capability, since idempotency, admissibility, and the front benchmark are
all statements about a single mesh.

\section{Summary}

The variational principle of \cite{mota2013} --- recover the internal
variable field by minimizing the $L^2$ distance to the source data in the
variable's own space --- already contains the answer for norm-like
variables. Their own space is a convex cone rather than a group, so the
minimization runs over the admissible convex subset of the finite element
space and becomes a variational inequality: a bound-constrained quadratic
program on the already-assembled Gram matrix, solved by a primal--dual
active-set method with a few additional back-substitutions.
Idempotency and stability under repeated remap follow from convexity
rather than requiring preservation; the logarithmic transformation
retains its role for asymptotic boundaries and composes with the
constraints where a variable possesses boundaries of both kinds; and the
active set of the constraint serves also as the front indicator that an
adaptive remeshing loop requires.

\begin{thebibliography}{9}

\bibitem{mota2013}
A.~Mota, W.~Sun, J.~T.~Ostien, J.~W.~Foulk~III, K.~N.~Long.
Lie-group interpolation and variational recovery for internal variables.
\emph{Computational Mechanics} 52 (2013) 1281--1299.

\bibitem{hintermuller2003}
M.~Hinterm\"uller, K.~Ito, K.~Kunisch.
The primal--dual active set strategy as a semismooth Newton method.
\emph{SIAM Journal on Optimization} 13 (2003) 865--888.

\bibitem{boyle1986}
J.~P.~Boyle, R.~L.~Dykstra.
A method for finding projections onto the intersection of convex sets in
Hilbert spaces.
In: \emph{Advances in Order Restricted Statistical Inference},
Lecture Notes in Statistics 37, Springer, 1986, 28--47.

\bibitem{lawson1974}
C.~L.~Lawson, R.~J.~Hanson.
\emph{Solving Least Squares Problems}.
Prentice-Hall, 1974.

\end{thebibliography}

\end{document}
